\QHeading{第1题　口算。}
\par\noindent


\begin{oralgrid}
\OralItem{1}{25\times4}{100}
\OralItem{2}{120\div6}{20}
\OralItem{3}{30\times8}{240}
\OralItem{4}{480\div60}{8}
\OralItem{5}{720\div8}{90}
\OralItem{6}{20\times30}{600}
\OralItem{7}{90\div30}{3}
\OralItem{8}{11\times7}{77}
\OralItem{9}{360\div9}{40}
\OralItem{10}{56\times10}{560}
\OralItem{11}{90\div5}{18}
\OralItem{12}{248\div8}{31}\end{oralgrid}
\par

\QHeading{口算思路补充（对应上表题号）}

\Answer{上面列出得数，下面按小题给出不同计算思路。计数单位法适用于小数，整数题可比较凑整、拆分或商不变等方法。}
\MethodTitle{第1小题：$25\times 4$}

\MethodBlock{口算思路一：按数位拆分计算}{把一个因数拆成整十、整百数与余下部分之和。}{\begin{aligned}[t]&\phantom{={}}25\times 4\\
&=4\times (20+5)\\
&=4\times 20+4\times 5\\
&=80+20\\
&=100\end{aligned}}
\MethodTitle{第2小题：$120\div 6$}

\MethodBlock{口算思路一：利用商不变的性质}{被除数和除数同时乘同一个非零数，商不变。}{\begin{aligned}[t]&\phantom{={}}120\div 6\\
&=(120\times 5)\div (6\times 5)\\
&=600\div 30\\
&=20\end{aligned}}

\MethodBlock{口算思路二：拆分除数，连续相除}{}{\begin{aligned}[t]&\phantom{={}}120\div 6\\
&=120\div (2\times 3)\\
&=120\div 2\div 3\\
&=60\div 3\\
&=20\end{aligned}}
\MethodTitle{第3小题：$30\times 8$}

\MethodBlock{口算思路一：凑整后补偿}{把接近整十、整百的因数凑整，再补回或减去多算的部分。}{\begin{aligned}[t]&\phantom{={}}30\times 8\\
&=30\times (10-2)\\
&=30\times 10-30\times 2\\
&=300-60\\
&=240\end{aligned}}
\MethodTitle{第4小题：$480\div 60$}

\MethodBlock{口算思路一：同时去掉相同个数的0}{}{\begin{aligned}[t]&\phantom{={}}480\div 60\\
&=(480\div 10)\div (60\div 10)\\
&=48\div 6\\
&=8\end{aligned}}

\MethodBlock{口算思路二：拆分除数，连续相除}{}{\begin{aligned}[t]&\phantom{={}}480\div 60\\
&=480\div (2\times 30)\\
&=480\div 2\div 30\\
&=240\div 30\\
&=8\end{aligned}}
\MethodTitle{第5小题：$720\div 8$}

\MethodBlock{口算思路一：利用商不变的性质}{被除数和除数同时乘同一个非零数，商不变。}{\begin{aligned}[t]&\phantom{={}}720\div 8\\
&=(720\times 5)\div (8\times 5)\\
&=3600\div 40\\
&=90\end{aligned}}

\MethodBlock{口算思路二：拆分除数，连续相除}{}{\begin{aligned}[t]&\phantom{={}}720\div 8\\
&=720\div (2\times 4)\\
&=720\div 2\div 4\\
&=360\div 4\\
&=90\end{aligned}}
\MethodTitle{第7小题：$90\div 30$}

\MethodBlock{口算思路一：同时去掉相同个数的0}{}{\begin{aligned}[t]&\phantom{={}}90\div 30\\
&=(90\div 10)\div (30\div 10)\\
&=9\div 3\\
&=3\end{aligned}}

\MethodBlock{口算思路二：拆分除数，连续相除}{}{\begin{aligned}[t]&\phantom{={}}90\div 30\\
&=90\div (2\times 15)\\
&=90\div 2\div 15\\
&=45\div 15\\
&=3\end{aligned}}
\MethodTitle{第8小题：$11\times 7$}

\MethodBlock{口算思路一：凑整后补偿}{把接近整十、整百的因数凑整，再补回或减去多算的部分。}{\begin{aligned}[t]&\phantom{={}}11\times 7\\
&=7\times (10+1)\\
&=7\times 10+7\times 1\\
&=70+7\\
&=77\end{aligned}}
\MethodTitle{第9小题：$360\div 9$}

\MethodBlock{口算思路一：拆分除数，连续相除}{}{\begin{aligned}[t]&\phantom{={}}360\div 9\\
&=360\div (3\times 3)\\
&=360\div 3\div 3\\
&=120\div 3\\
&=40\end{aligned}}
\MethodTitle{第10小题：$56\times 10$}

\MethodBlock{口算思路一：按数位拆分计算}{把一个因数拆成整十、整百数与余下部分之和。}{\begin{aligned}[t]&\phantom{={}}56\times 10\\
&=10\times (50+6)\\
&=10\times 50+10\times 6\\
&=500+60\\
&=560\end{aligned}}
\MethodTitle{第11小题：$90\div 5$}

\MethodBlock{口算思路一：拆分被除数}{}{\begin{aligned}[t]&\phantom{={}}90\div 5\\
&=(50+40)\div 5\\
&=50\div 5+40\div 5\\
&=10+8\\
&=18\end{aligned}}

\MethodBlock{口算思路二：利用商不变的性质}{被除数和除数同时乘同一个非零数，商不变。}{\begin{aligned}[t]&\phantom{={}}90\div 5\\
&=(90\times 2)\div (5\times 2)\\
&=180\div 10\\
&=18\end{aligned}}
\MethodTitle{第12小题：$248\div 8$}

\MethodBlock{口算思路一：拆分被除数}{}{\begin{aligned}[t]&\phantom{={}}248\div 8\\
&=(240+8)\div 8\\
&=240\div 8+8\div 8\\
&=30+1\\
&=31\end{aligned}}

\MethodBlock{口算思路二：利用商不变的性质}{被除数和除数同时乘同一个非零数，商不变。}{\begin{aligned}[t]&\phantom{={}}248\div 8\\
&=(248\times 5)\div (8\times 5)\\
&=1240\div 40\\
&=31\end{aligned}}


\par\noindent

\QHeading{第2题　脱式计算。}

\CalcBlock{1}{\begin{aligned}[t]&\phantom{={}}\color{black}12\times3+45\div9\\&\color{answerred}=36+5\\&\color{answerred}=41\end{aligned}}
\par

\CalcBlock{2}{\begin{aligned}[t]&\phantom{={}}\color{black}25\times4-72\div8\\&\color{answerred}=100-9\\&\color{answerred}=91\end{aligned}}
\CalcBlock{3}{\begin{aligned}[t]&\phantom{={}}\color{black}15\times6+30\times7\\&\color{answerred}=90+210\\&\color{answerred}=300\end{aligned}}
\CalcBlock{4}{\begin{aligned}[t]&\phantom{={}}\color{black}45\times2-18\times5\\&\color{answerred}=90-90\\&\color{answerred}=0\end{aligned}}
\CalcBlock{5}{\begin{aligned}[t]&\phantom{={}}\color{black}96\div6+63\div3\\&\color{answerred}=16+21\\&\color{answerred}=37\end{aligned}}
\CalcBlock{6}{\begin{aligned}[t]&\phantom{={}}\color{black}1000\div5-24\times2\\&\color{answerred}=200-48\\&\color{answerred}=152\end{aligned}}
\par\noindent

\QHeading{第3题　解决问题。}

\Problem{①李老师带2000元钱去买排球和篮球，每个排球50元，每个篮球60元。李老师买20个排球和15个篮球，一共花了多少钱？}
\par


\MethodTitle{方法一}
\Answer{\Formula{\begin{aligned}[t]&\phantom{={}}50\times20+60\times15\\&=1000+900\\&=1900\text{（元）}\end{aligned}}\par 答：一共花了1900元。}

\MethodTitle{方法二：先统一单价，再补差价}
\Answer{先把篮球也按每个50元计算，再补上每个篮球少算的10元。}
\Answer{\Formula{\begin{aligned}[t]&\phantom{={}}50\times (20+15)+(60-50)\times 15\\
&=50\times 35+10\times 15\\
&=1750+150\\
&=1900\text{（元）}\end{aligned}}}
\Answer{答：一共花了1900元。}
\Problem{②春光小学四年级有6个班，平均每班有学生45人；三年级有8个班，平均每班有学生42人。四年级比三年级少多少人？}

\MethodTitle{方法一}
\Answer{\Formula{\begin{aligned}[t]&\phantom{={}}42\times8-45\times6\\&=336-270\\&=66\text{（人）}\end{aligned}}\par 答：四年级比三年级少66人。}

\MethodTitle{方法二：按相同班数比较}
\Answer{三年级多2个班，按每班42人多84人；相同的6个班中，四年级每班又多3人，要扣去这部分。}
\Answer{\Formula{\begin{aligned}[t]&\phantom{={}}42\times (8-6)-(45-42)\times 6\\
&=42\times 2-3\times 6\\
&=84-18\\
&=66\text{（人）}\end{aligned}}}
\Answer{答：四年级比三年级少66人。}
